\begin{center}
    {\LARGE \textbf{2016分析化学}} \\[2ex]
    {\large 时量180分钟 \quad 满分90分}
\end{center}

\vspace{0.5cm}
\noindent\textbf{注意事项:}
\begin{enumerate}[label=\arabic*., parsep=0pt, itemsep=0pt]
    \item 答题前,考生先将自己的姓名、学号、考场号、座位号填写在答题卡上,将条形码准确粘贴在考生信息条形码粘贴区。
    \item 考试结束后,将本试卷与答题纸一并收回。
    \item 回答各个问题时,将答案写在答题纸上,写在本试卷上无效。
\end{enumerate}

\vspace{0.5cm}
\noindent\textbf{一、简答题(42 pts.)}

\begin{enumerate}[label=\arabic*., itemsep=1.5ex]
    \item 在$\ce{K2Cr2O7}$标定$\ce{Na2S2O3}$的实验中，为何不采用直接法标定，而是采用间接碘量法标定？
    \item 写出下列溶液的质子条件式
    \begin{enumerate}[label=(\arabic*), noitemsep, topsep=0pt, partopsep=0pt]
        \item $60\ \mathrm{mL}\ 0.10\ \mathrm{mol/L}\ \ce{Na2CO3}$与$40\ \mathrm{mL}\ 0.15\ \mathrm{mol/L}\ \ce{HCl}$的混合溶液
        \item $20\ \mathrm{mL}\ 0.01\ \mathrm{mol/L}\ \ce{NaOH}$和$10\ \mathrm{mL}\ 0.10\ \mathrm{mol/L}\ \ce{H2SO4}$的混合溶液
    \end{enumerate}
    \item 根据下表所给数据，计算用$0.010\ \mathrm{mol/L}\ \ce{Ce^{4+}}$滴定同浓度$\ce{Fe^{2+}}$时表中各点$\varphi$值（$\mathrm{V}$）
    \begin{center}
    \begin{tabular}{|c|c|c|c|}
    \hline
    浓度 & 计量点前$0.1\%$ & 化学计量点 & 计量点后$0.1\%$ \\
    \hline
    $0.10\ \mathrm{mol/L}$ & $0.86$ & $1.06$ & $1.26$ \\
    \hline
    $0.010\ \mathrm{mol/L}$ &  &  &  \\
    \hline
    \end{tabular}
    \end{center}
    \item 用佛尔哈德返滴定法测定$\ce{Cl^-}$时，为何在形成$\ce{AgCl}$沉淀后向溶液中加入硝基苯，若不加硝基苯对测定结果有何影响？
    \item 测定黄铁矿中硫的质量分数，六次测定结果分别为$30.48\%$，$30.42\%$，$30.59\%$，$30.51\%$，$30.56\%$，$30.49\%$，计算置信水平为$95\%$时总体平均值的置信区间。已知置信度为$95\%$时，$t_{0.05,5}=2.57$
    \item 计算$\ce{Hg^{2+}/Hg2^{2+}}$电对在$[\ce{CN^-}]=0.1\ \mathrm{mol/L}$溶液中的条件电位(忽略离子强度的影响)。已知$\varphi^\ominus_{\ce{Hg^{2+}/Hg2^{2+}}}=0.907\ \mathrm{V}$，$K_{\mathrm{sp}}(\ce{Hg2(CN)2})=5.0\times10^{-40}$，$\ce{Hg-CN^-}$配合物的$\lg\beta_4=41.40$
    \item 已知$\ce{Zn^{2+}}$在二苯硫腙$-CHCl_3$和水中的分配比为$80$，若以$50\ \mathrm{mL}$萃取剂分两次萃取含有$0.015\ \mathrm{g}\ \ce{Zn^{2+}}$的$100\ \mathrm{mL}$水溶液，问水相中还有多少克$\ce{Zn^{2+}}$？
\end{enumerate}

\vspace{0.5cm}
\noindent\textbf{二、计算题(48 pts.)}

\begin{enumerate}[label=\arabic*., itemsep=1.5ex]
    \item 在$\mathrm{pH}=10.0$的氨性缓冲液中，以铬黑$\mathrm{T(EBT)}$作指示剂，用$0.020\ \mathrm{mol/L}\ \mathrm{EDTA}$滴定$0.020\ \mathrm{mol/L}$的$\ce{Zn^{2+}}$，终点时游离氨的浓度为$0.20\ \mathrm{mol/L}$，若终点时$\mathrm{pZn_{ep}}=12.2$，计算终点误差。已知$\lg K_{\ce{ZnY}}=16.5$，$\mathrm{pH}=10.0$时$\lg\alpha_{Y(\mathrm{H})}=0.45$，$\lg\alpha_{\ce{Zn(OH)}}=2.4$，$\ce{Zn-NH3}$的$\lg\beta_1-\lg\beta_4$分别为$2.37$，$4.81$，$7.31$，$9.46$
    \item 计算$\ce{CaC2O4}$在$\mathrm{pH}=5.0$，草酸总浓度为$0.05\ \mathrm{mol/L}$溶液中的溶解度。已知$\ce{CaC2O4}$的$\mathrm{p}K_{\mathrm{sp}}=7.8$，$\ce{H2C2O4}$的$\mathrm{p}K_{a1}=1.22$，$\mathrm{p}K_{a2}=4.19$
    \item 用$0.10\ \mathrm{mol/L}\ \ce{NaOH}$溶液滴定同浓度邻苯二甲酸氢钾（简写成$\ce{KHP}$）。计算化学计量点及其前后$0.1\%$的$\mathrm{pH}$值。已知$\ce{H2P}$的$\mathrm{p}K_{a1}=2.95$，$\mathrm{p}K_{a2}=5.41$
    \item $\ce{Ti}$和$\ce{V}$与$\ce{H2O2}$作用生成有色络合物，分别各称取$5.00\ \mathrm{mg}$的纯金属均用$\ce{HClO4}$及$\ce{H2O2}$处理后都稀释至$100\ \mathrm{mL}$作为标准溶液。另取$0.800\ \mathrm{g}$含$\ce{Ti}$和$\ce{V}$的合金用同样方法处理。这三份溶液各用厚度为$1\ \mathrm{cm}$的吸收皿在$410\ \mathrm{nm}$及$460\ \mathrm{nm}$处测得吸光度值如下：
    \begin{center}
    \begin{tabular}{|c|c|c|}
    \hline
    溶液 & $\mathrm{A\ (410nm)}$ & $\mathrm{A\ (460nm)}$ \\
    \hline
    $\ce{Ti}$ & $0.760$ & $0.513$ \\
    \hline
    $\ce{V}$ & $0.185$ & $0.250$ \\
    \hline
    合金 & $0.800$ & $0.690$ \\
    \hline
    \end{tabular}
    \end{center}
    计算合金中$\ce{Ti}$、$\ce{V}$的含量。已知$A_r(\ce{Ti})=47.88$，$A_r(\ce{V})=50.94$
\end{enumerate}